\scanpage{38}\mspage{17}

alors l'image $\bar M$ de $M$ dans $\unsure{\mathrm{Ep}|\mathrm{E}p}(\deleted{M_0}E_0)$ ($\deleted{M_0}E_0=$ ens. des
éléments minimaux de $M$) est un groupe. Donc
si $\inserted{Y=}B_M$ est $1$-connexe, ce groupe est réduit à $\{1\}$,
ce qui signifie que pour $\deleted{\unsure{\psi}}\lambda\in\deleted{M_0}E_0$ (p. ex.
pour $\lambda=\psi$), \deleted{et $u$} et $u\in M$, on a $u\lambda=\lambda$,
ce qui implique que $M$ est cofiltrant.

Dans le cas contraire $(\varphi^r\psi=\varphi^s\psi\Longrightarrow r=s)$
prenant la relation d'équivalence \inserted{inv. à droite} $R_\varphi$ \inserted{sur $M$} définie par
$u\overset{R_\varphi}{\sim}v\Longleftrightarrow\exists p,q\in\mathbb N$ avec $\varphi^p u=\varphi^q v$, et
le quotient $F=M/R_\varphi$, \unsure{vu|considéré} comme $M$-ens.
à droite, i.e. comme objet de $Y^\wedge$, alors
$Y_{/F}$ est \inserted{(pseudocofiltrant et)} $0$-connexe (tous ses objets sont $\unsure{=}$
équivalents), mais \deleted{\unsure{il|elle}} n'est pas $1$-connexe ;
plus précisément son groupe fondamental
est isomorphe à $\pi_1(M_0)$, où \unsure{1}
\[
M_0=\{f\in M\mid\psi f\overset{R_\varphi}{\sim}\psi\},
\]
et admet un groupe quotient isomorphe à $\mathbb Z$,
grâce à l'hom. surjectif
\diagram{\begin{tikzcd}[ampersand replacement=\&]
\pi_1(M_0)=G_0\arrow[r,"\delta_0"] \& \mathbb Z
\end{tikzcd}}
donné sur $M_0$ par
\[
\delta(f)=q-p\qquad\text{si}\qquad\varphi^p\underline{\psi f}=\varphi^q\psi.
\]

[ \uline{NB} \deleted{On} \deleted{\unsure{Il exi}} Comme $\psi\in M_0$, on a $\varphi\psi\in M_0$,
\deleted{\unsure{de même}} donc $\exists f$ tel que $\psi f=\varphi\psi$,
et alors $f\in M_0$ et $\delta(f)=1$ ] Donc on a $\unsure{=}$
$H^{1}(X,\unsure{\mathbb Q|Q|\mathbb Z})\ne0$.

% archive sheet number: 37
